% TOP_PROBLEM: {"id": 1200009, "title": "Some Questions — Question 9", "category_id": 17, "category": {"id": 17, "name": "group_theory", "display_name": "Group Theory", "description": "Problems about groups, group actions, representations, and related algebraic structures.", "slug": "group-theory", "order_index": 17, "created_at": "2026-07-31T15:26:25.671Z"}, "status": "open", "problem_number": "AMR-011-0009", "set_id": 13}
% FIRST_POSTED: 2026-10-01
% ENTRY_KIND: statement_only
% UnsolvedMath ID 1200009; source catalogue code AMR-011-0009
% !TeX program = lualatex
% Definition and sources only; this file is not an LLM solution attempt.
\documentclass[11pt,a4paper]{article}
\usepackage[margin=25mm]{geometry}
\usepackage{fontspec}
\setmainfont{lmroman10-regular.otf}[BoldFont=lmroman10-bold.otf,
 ItalicFont=lmroman10-italic.otf,BoldItalicFont=lmroman10-bolditalic.otf]
\usepackage{amsmath,amssymb,mathtools}
\usepackage{unicode-math}
\setmathfont{latinmodern-math.otf}
\AtBeginDocument{\let\setminus\smallsetminus}
\usepackage{xurl,hyperref}
\hypersetup{colorlinks=true,urlcolor=blue}
\setcounter{secnumdepth}{0}
\setlength{\parindent}{0pt}
\setlength{\parskip}{5pt}
\setlength{\emergencystretch}{3em}
\begin{document}
\section*{Some Questions — Question 9}\label{1200009}
{\small UnsolvedMath ID 1200009 · AMR-011-0009 · Group Theory · AMR Open Problem Lists}
\subsection{Definitions and mathematical statement}
Let $\Gamma$ be generated by a finite symmetric set $S$, and let $H_1,H_2,\ldots$ be finite-index normal subgroups. For each quotient $Q_n=\Gamma/H_n$, use the Cayley multigraph defined by the images of $S$, retaining generator multiplicities and loops if they occur. Its normalized averaging operator is
\[
 P_n f(q)=|S|^{-1}\sum_{s\in S}f(qs).
\]
The family has property $\tau$ if some $\varepsilon>0$ satisfies
$\langle(I-P_n)f,f\rangle\geq\varepsilon\langle f,f\rangle$ for every $n$ and every real function $f$ of mean zero on $Q_n$. Inner products use uniform probability measure. This is the upper spectral-gap convention, not an exclusion of bipartiteness.

Set $\Gamma_n=\bigcap_{j=1}^nH_j$, a descending chain of finite-index normal subgroups. If the original family has property $\tau$, must the quotient family $\Gamma/\Gamma_n$ also have property $\tau$?

The new positive gap may differ from the original one but must be independent of $n$. Neither family is assumed to have trivial total intersection, and the original subgroups need not be nested. Although each new quotient maps onto all the earlier quotients, a gap for those images does not by definition control every function on the larger quotient. The question asks for the stability of uniform expansion under these successive finite intersections, not merely the existence of a spectral gap in each individual finite connected quotient.
\subsection{Short English statement}
Does property tau survive passing from a family of normal finite-index subgroups to its successive intersections?
\subsection{Sources}
\begin{itemize}
\item[S1] ULAM.AI and the UnsolvedMath Contributors, supplied problems.json, record 1200009 (AMR-011-0009), AMR Open Problem Lists. Catalogue identity and source transcription; CC BY 4.0.
\url{https://huggingface.co/datasets/ulamai/UnsolvedMath}.
\item[S2] Abert - Some questions (pdf) (2010), Question 9. Original source indexed by AMR-011-0009.
\url{https://www.renyi.hu/~abert/questions.pdf}.
\end{itemize}
\end{document}
